\documentclass[dvipsnames]{beamer}
\usepackage{comment}
\usepackage{hyperref}

\title{MATH 1190 \& 1210 Meeting for Week 12}
\author{Josh}
\date{30 March 2026}

\newcommand{\transition}[1]{\frame{{}\begin{center}\fbox{#1}\end{center}}}

\begin{document}

\maketitle


\section{Logistics: Pacing}

\transition{Logistics}

\frame{{Logistics: Pacing}

{\bf Where are folks in the lessons right now?}

}

\frame[t]{{Logistics: Upcoming Schedule}

    \begin{itemize}
      \item week 12 of class (30 March - 3 April)
        \begin{itemize}
        \item lesson 18 (FTOC) 
        \item lesson 19 (Summations)
        \item KC 3 on Thursday 2 April 
        \end{itemize}
     \item week 13 of class (6 April - 10 April)
        \begin{itemize}
        \item lesson 20 (Discrete RVs)
        \item Application Day 3: Gradient Descent
        \end{itemize}
    \end{itemize}

}

\frame[t]{{Logistics: Upcoming Deadlines}

{\bf Assessment-related Logistics}
\begin{itemize}
\item New Pre-Class for lesson 18 (FTOC)
\item The KC 3 grading
    \begin{itemize}
        \item {\color{RedOrange}A2}:
        \item {\color{RedOrange}A3}:
        \item {\color{Fuchsia}I1}:
        \item {\color{Fuchsia}I2}:
    \end{itemize}




Reflection is due on Tuesday 24 March before 11:59 PM. Good opportunity to give students feeback and promote metacognition.
\item Assessment 2 regrade requests open now until Thursday 26 March.
\item Assessment 3
    \begin{itemize}
    \item new targets: {\color{RedOrange}A2},{\color{RedOrange}A3}, {\color{Fuchsia}I1}, {\color{Fuchsia}I2}, 
    \item reassessment targets: {\color{ProcessBlue}F2}, {\color{ForestGreen}D2}, {\color{ForestGreen}D3}, {\color{RedOrange}A1}
    \item Tuesday, 31 March, 7:00 PM; 80 minutes
    \item room assignments are currently the same as KC1 unless folks feel like they won't be ready and we need to change the date.
    \end{itemize}
\end{itemize}

}




\section{Content}

\transition{Content}


%------------------------------------------------

\begin{frame}{Week Overview}
\textbf{Main Topics:}
\begin{itemize}
\item Antiderivatives of basic functions (I3)
\item Computing definite integrals using FTC (I4)
\end{itemize}

\vspace{1em}

\textbf{Core Goals:}
\begin{itemize}
\item Build fluency with basic antiderivatives
\item Connect antiderivatives to definite integrals
\item Understand when FTC applies (and when it does not)
\end{itemize}
\end{frame}

%------------------------------------------------

\begin{frame}{Learning Objectives}
Students should be able to:
\begin{itemize}
\item Compute basic antiderivatives
\item Use the Fundamental Theorem of Calculus (FTC)
\item Check continuity of the integrand
\item Simplify expressions before integrating
\item Interpret definite integrals as net area
\end{itemize}
\end{frame}

%------------------------------------------------

\begin{frame}{Lesson Flow}
\begin{itemize}
\item Warm-up: basic indefinite integrals
\item Mini-lecture: Fundamental Theorem of Calculus
\item Exercise A: When does FTC apply?
\item Example 1: Algebra + FTC
\item Exercise B: Full integration workflow
\end{itemize}

\vspace{1em}

\end{frame}

%------------------------------------------------

\begin{frame}{Warm-Up: Antiderivatives}
\textbf{Purpose:}
\begin{itemize}
\item Reinforce basic rules (power, exponential)
\item Prepare students for FTC problems
\end{itemize}

\vspace{1em}

\textbf{Instructor Notes:}
\begin{itemize}
\item Encourage students to show steps
\item Watch for common errors with exponents
\item Reinforce the idea of a \textbf{family of functions}
\end{itemize}
\end{frame}

%------------------------------------------------

\begin{frame}{Mini-Lecture: Fundamental Theorem of Calculus}
\textbf{Big Idea:}
\begin{itemize}
\item Definite integrals connect to antiderivatives
\end{itemize}

\vspace{1em}

\textbf{Key Message:}
\begin{itemize}
\item $\int_a^b f(x),dx = F(b) - F(a)$ \textbf{only if $f$ is continuous}
\end{itemize}

\vspace{1em}

\textbf{Important nuance:}
\begin{itemize}
\item If $f$ is not continuous:
\begin{itemize}
\item FTC may not apply
\item Integral may still exist
\end{itemize}
\end{itemize}
\end{frame}

%------------------------------------------------



%------------------------------------------------

\begin{frame}{Exercise A: Does FTC Apply?}
\textbf{Goal:}
\begin{itemize}
\item Identify whether integrand is continuous on $[a,b]$
\end{itemize}

\vspace{1em}

\textbf{Common Issues:}
\begin{itemize}
\item Students skip checking continuity
\item Confusion with removable vs non-removable discontinuities
\end{itemize}

\vspace{1em}

\textbf{Teaching tip:}
\begin{itemize}
\item Make this a required step before applying FTC
\end{itemize}
\end{frame}

%------------------------------------------------

\begin{frame}{Example 1: Algebra + FTC}
\textbf{Purpose:}
\begin{itemize}
\item Show how algebra simplifies integration
\end{itemize}

\vspace{1em}

\textbf{Key Skills:}
\begin{itemize}
\item Simplify integrand first
\item Then apply basic antiderivatives
\item Then apply FTC
\end{itemize}

\vspace{1em}

\textbf{Instructor choice:}
\begin{itemize}
\item Model solution OR scaffold student work
\end{itemize}
\end{frame}

%------------------------------------------------

\begin{frame}{Exercise B: Full Workflow}
\textbf{Students must:}
\begin{itemize}
\item Simplify the integrand
\item Find antiderivatives
\item Check FTC conditions
\item Evaluate definite integral
\end{itemize}

\vspace{1em}
\end{frame}
%---------------------------------------------

\begin{frame}{Common Student Difficulties}
\begin{itemize}
\item Skipping algebra simplification
\item Forgetting to check continuity
\item Misusing FTC when conditions fail
\item Treating definite integrals as purely procedural
\end{itemize}
\end{frame}

%------------------------------------------------

\begin{frame}{Net Area vs Area}
\textbf{Key distinction:}
\begin{itemize}
\item Definite integral = \textbf{net (signed) area}
\item Area = always positive
\end{itemize}

\vspace{1em}

\textbf{Use Extra Practice 2 if time allows}
\begin{itemize}
\item Helps clarify conceptual understanding
\end{itemize}
\end{frame}

%------------------------------------------------

\begin{frame}{Discussion for Instructors}
\textbf{Where students may struggle:}
\begin{itemize}
\item Understanding FTC conceptually
\item Knowing when it applies
\item Connecting algebra to integration
\end{itemize}

\vspace{1em}

\textbf{Questions:}
\begin{itemize}
\item What misconceptions have you seen before?
\item How do you emphasize the FTC hypothesis?
\end{itemize}
\end{frame}

%------------------------------------------------


%------------------------------------------------

\begin{frame}{Final Takeaway}
\textbf{Goal of this lesson:}
\begin{itemize}
\item Build a coherent process:
\begin{itemize}
\item Simplify → Antiderivative → Check → Evaluate
\end{itemize}
\end{itemize}

\vspace{1em}

\textbf{Big idea:}
\begin{itemize}
\item FTC is not just a formula—it’s a \textbf{conceptual bridge}
\end{itemize}
\end{frame}

%------------------------------------------------

%%%%%%%%%%%%%%%%%%%%%%%%%%%%%%%%%%%%%%%%%%%%%%%%%%%%%%%%%%%%%%%%%%%%%%%%%%%%%%%%%%%%%%%%%%%%%%%%%%%%%%
%%%%%%%%%%%%%%%%%%%%%%%%%%%%%%%%%%%%%%%%%%%%%%%%%%%%%%%%%%%%%%%%%%%%%%%%%%%%%%%%%%%%%%%%%%%%%%%%%%%%%%




%------------------------------------------------

\begin{frame}{Week Overview}
\textbf{Main Topic: Summation Notation (F3)}

\vspace{1em}

Students will learn to:
\begin{itemize}
\item Compute sums from sigma notation
\item Write expressions using summation notation
\item Apply linearity properties of sums
\end{itemize}

\vspace{1em}

\textbf{Instructor Focus:}
\begin{itemize}
\item Emphasize \textbf{process over speed}
\item Encourage careful handling of indices
\item Reinforce pattern recognition
\end{itemize}
\end{frame}

%------------------------------------------------

\begin{frame}{Exercise Set 1: Evaluating Sums}
\textbf{Exercises: 1(a), 1(b), 1(c)}

\vspace{1em}

\textbf{Core Skills:}
\begin{itemize}
\item Substituting index values correctly
\item Evaluating expressions step-by-step
\item Handling fractions, logs, and alternating signs
\end{itemize}

\vspace{1em}

\textbf{Common Student Issues:}
\begin{itemize}
\item Skipping steps → arithmetic mistakes
\item Confusion about starting index
\item Errors with negative signs in alternating sums
\end{itemize}
\end{frame}

%------------------------------------------------

\begin{frame}{Instructor Notes: Evaluating Sums}
\textbf{What to emphasize:}
\begin{itemize}
\item Write out \textbf{at least the first few terms}
\item Encourage students to check patterns
\item Slow down when alternating signs appear
\end{itemize}

\vspace{1em}

\textbf{Key teaching point:}
\begin{itemize}
\item Students should see summation as a \textbf{process}, not a shortcut
\end{itemize}
\end{frame}

%------------------------------------------------

\begin{frame}{Exercise Set 2: Writing Summations}
\textbf{Exercises: 2(a), 2(b), 2(c)}

\vspace{1em}

\textbf{Core Skills:}
\begin{itemize}
\item Identifying patterns in sequences
\item Choosing appropriate indices
\item Writing compact sigma notation
\end{itemize}

\vspace{1em}

\textbf{Key patterns:}
\begin{itemize}
\item Even numbers → $2k$
\item Alternating signs → $(-1)^k$
\item Powers → $a^k$
\end{itemize}
\end{frame}

%------------------------------------------------

\begin{frame}
\textbf{What to emphasize:}
\begin{itemize}
\item Help students identify:
\begin{itemize}
\item Numerator pattern
\item Denominator pattern
\item Sign pattern
\end{itemize}
\end{itemize}

\vspace{1em}

\textbf{Common mistake:}
\begin{itemize}
\item Correct pattern, but incorrect index bounds
\end{itemize}
\end{frame}

%------------------------------------------------

\begin{frame}{Exercise Set 3: Linearity of Sums}
\textbf{Exercises: 3(a), 3(b)}

\vspace{1em}

\textbf{Core Skills:}
\begin{itemize}
\item Splitting sums into parts
\item Factoring constants
\item Keeping index-dependent terms inside sums
\end{itemize}

\vspace{1em}

\textbf{Big Idea:}
\begin{itemize}
\item Summation behaves like algebra—\textbf{but not always}
\end{itemize}
\end{frame}

%------------------------------------------------

\begin{frame}{Linearity}
\textbf{Emphasize the difference:}

\vspace{1em}

\textbf{Valid:}
\begin{itemize}
\item $\sum (2a_k) = 2\sum a_k$
\end{itemize}

\textbf{Be careful:}
\begin{itemize}
\item Terms involving $k$ cannot be factored out
\end{itemize}

\vspace{1em}

\textbf{Common issue:}
\begin{itemize}
\item Students over-generalize factoring rules
\end{itemize}
\end{frame}

%------------------------------------------------

\begin{frame}{Connecting to Learning Goals (F3)}
\textbf{Students should:}
\begin{itemize}
\item Compute sums accurately
\item Translate expressions into summation notation
\item Apply linearity correctly
\end{itemize}

\end{frame}

%------------------------------------------------

\begin{frame}{Discussion}
\textbf{Where students may struggle:}
\begin{itemize}
\item Indexing and bounds
\item Recognizing patterns
\item Misapplying algebra rules
\end{itemize}

\vspace{1em}

\textbf{Questions:}
\begin{itemize}
\item Which of these issues have you seen before?
\item What strategies help students most here?
\end{itemize}
\end{frame}



%------------------------------------------------
%%%%%%%%%%%%%%%%%%%%%%%%%%%%%%%%%%%%%%%%%%%%%%%%%%%%%%%%%%%%%%%%%%%%%%%%%%%%%%%%%%%%%%%%%%%%%%%%%%%%%%
%%%%%%%%%%%%%%%%%%%%%%%%%%%%%%%%%%%%%%%%%%%%%%%%%%%%%%%%%%%%%%%%%%%%%%%%%%%%%%%%%%%%%%%%%%%%%%%%%%%%%%



%%%%%%%%%%%%%%%%%%%%%%%%%%%%%%%%%%%%%%%%%%%%%%%%%%%%%%%%%%%%%%%%%%%%%%%%%%%%%%%%%%%%%%%%%%%%%%%%%%%%%%
%%%%%%%%%%%%%%%%%%%%%%%%%%%%%%%%%%%%%%%%%%%%%%%%%%%%%%%%%%%%%%%%%%%%%%%%%%%%%%%%%%%%%%%%%%%%%%%%%%%%%%

\section{Pedagogy}

\transition{Pedagogy}


\transition{Active Learning}
\transition{And Now - An Active Learning Discussion with Jim B.!}
%%%%%%%%%%%%%%%%%%%%%%%%%%%%%%%%%%%%%%%%%%%%%%%%%%%%%%%%%%%%%%%%%%%%%%%%%%%%%%%%%%%%%%%%%%%%%%%%%%%%%%

%------------------------------------------------

\begin{frame}{Reading Focus for Today}
\textbf{Please read:}
\begin{itemize}
\item Section 3.1: \textit{Educational Practice: Assessment and Grading}
\end{itemize}

\vspace{1em}

\textbf{As you read, look for:}
\begin{itemize}
\item What messages do grades send to students?
\item What does “diagnostic feedback” look like?
\item How does assessment shape student mindset?
\end{itemize}
\end{frame}

%------------------------------------------------

\begin{frame}{Key Idea: Performance Culture vs Learning Culture}
\textbf{Performance Culture}
\begin{itemize}
\item Focus on scores and grades
\item Students compare themselves to others
\item Little guidance on how to improve
\end{itemize}

\vspace{1em}

\textbf{Assessment for Learning (A4L)}
\begin{itemize}
\item Feedback replaces (or supplements) grades
\item Focus on progress and improvement
\item Learning is a \textit{journey}
\end{itemize}
\end{frame}

%------------------------------------------------

\begin{frame}{Discussion 1}
\textbf{Reflect on our current courses:}
\begin{itemize}
\item Where do we see a “performance culture”?
\item What messages do our exams send?
\item Do students know how to improve after an exam?
\end{itemize}
\end{frame}

%------------------------------------------------

\begin{frame}{What is Effective Feedback?}
A4L suggests feedback should answer:
\begin{itemize}
\item Where is the student now?
\item Where do they need to go?
\item How can they close the gap?
\end{itemize}

\vspace{1em}

\textbf{Key Practice: Diagnostic Feedback}
\begin{itemize}
\item Identifies strengths
\item Identifies specific areas for improvement
\item Gives actionable next steps
\end{itemize}
\end{frame}

%------------------------------------------------

\begin{frame}{Research Insight: Butler’s Studies}
Three conditions:
\begin{itemize}
\item Grades only
\item Grades + comments
\item Comments only
\end{itemize}

\vspace{1em}

\textbf{Result:}
\begin{itemize}
\item \textbf{Comments only} → highest achievement & motivation
\item Grades reduce engagement (even with comments)
\end{itemize}
\end{frame}

%------------------------------------------------

\begin{frame}{Discussion 2}
\textbf{Implications for our exams:}
\begin{itemize}
\item Do students focus more on 4's and 5's or feedback?
\item How might we emphasize comments over scores?
\end{itemize}
\end{frame}

%------------------------------------------------

\begin{frame}{Rubrics and Self-Assessment}
Effective A4L classrooms use:
\begin{itemize}
\item Clear rubrics
\item Student self-assessment
\item Opportunities for revision
\end{itemize}

\vspace{1em}

\textbf{Key finding:}
\begin{itemize}
\item Students improve most when they track their own progress
\end{itemize}
\end{frame}

%------------------------------------------------

\begin{frame}{Discussion 3}
\textbf{Think about implementation:}
\begin{itemize}
\item Do we provide clear criteria for success?
\item Could students use our rubrics to self-assess?
\item Where could we add additional opportunities for revision?
\end{itemize}
\end{frame}





%------------------------------------------------
%%%%%%%%%%%%%%%%%%%%%%%%%%%%%%%%%%%%%%%%%%%%%%%%%%%%%%%%%%%%%%%%%%%%%%%%%%%%%%%%%%%%%%%%%%%%%%%%%%%%%%%%%%%%%%%%%%%%%%%%%%%%%%%%%%%%%%%%%%%%%%%%%%%%%%%%%%%%%%%%%%%%%%%%%%%%%%%%%%%%%%%%%%%%%%%%%%%%%%%%%%%%%%%%%%%%%%%%%%%%
\begin{frame}{Context: We Already Use Growth-Based Grading}
\textbf{Important starting point:}
\begin{itemize}
\item Our system is already designed around growth and mastery
\item Students often have opportunities to revise or reattempt
\end{itemize}

\vspace{1em}

\textbf{Key question:}
\begin{itemize}
\item Are our \textbf{feedback practices} aligned with this system?
\end{itemize}

\vspace{1em}

\textbf{Potential tension:}
\begin{itemize}
\item Growth-based grading system
\item But \textit{performance-based feedback} (points, correctness)
\end{itemize}
\end{frame}

%------------------------------------------------

\begin{frame}{Where Growth Systems Break Down}
Even in growth-based systems, problems arise when:
\begin{itemize}
\item Feedback is minimal (“-2 points”)
\item Comments are vague (“review chain rule”)
\item Students focus only on scores or levels
\item Revision is allowed but not well-supported
\end{itemize}

\vspace{1em}

\textbf{Result:}
\begin{itemize}
\item Students do not actually learn from feedback
\item System behaves like traditional grading
\end{itemize}
\end{frame}

%------------------------------------------------

\begin{frame}{Discussion: Our Current Practice}
\textbf{Be honest about our courses:}
\begin{itemize}
\item What do students actually read: comments or scores?
\item How detailed is our feedback on exams?
\item Do students know what to do differently next time?
\end{itemize}

\vspace{1em}

\textbf{Follow-up:}
\begin{itemize}
\item Where are we already doing this well?
\end{itemize}
\end{frame}

%------------------------------------------------

\begin{frame}{From “Correctness” to “Progress”}
\textbf{Traditional feedback:}
\begin{itemize}
\item Focus on right vs wrong
\item Emphasis on final answer
\end{itemize}

\vspace{1em}

\textbf{Growth-aligned feedback:}
\begin{itemize}
\item Focus on \textbf{process and reasoning}
\item Identify what was done correctly
\item Highlight the next step for improvement
\end{itemize}

\vspace{1em}

\textbf{Shift:}
\begin{itemize}
\item From evaluation → to guidance
\end{itemize}
\end{frame}

%------------------------------------------------


%------------------------------------------------







%------------------------------------------------

\begin{frame}{Example: Calculus Feedback}
\textbf{Less effective:}
\begin{itemize}
\item “Incorrect derivative”
\item “Algebra mistake”
\end{itemize}

\vspace{1em}

\textbf{More effective:}
\begin{itemize}
\item “You correctly applied the product rule setup.”
\item “Error occurs when differentiating the second term.”
\item “Revisit derivative of $e^x$ and retry from this step.”
\end{itemize}

\vspace{1em}

\textbf{Key elements:}
\begin{itemize}
\item Acknowledge progress
\item Locate the error
\item Give a next action
\end{itemize}
\end{frame}

%------------------------------------------------

\begin{frame}{Supporting Revision (Critical Piece)}
Research shows:
\begin{itemize}
\item Feedback only works if students \textbf{use it}
\end{itemize}

\vspace{1em}

\textbf{Questions for our system:}
\begin{itemize}
\item Do students actually revise their work?
\item Do we require reflection on mistakes?
\item Is feedback detailed enough to support revision?
\end{itemize}
\end{frame}

%------------------------------------------------

\begin{frame}{Discussion: Revision Practices}
\textbf{Compare approaches:}
\begin{itemize}
\item Full exam corrections?
\item Targeted skill reassessment?
\item Written reflection required?
\end{itemize}

\vspace{1em}

\textbf{Key question:}
\begin{itemize}
\item What revision structure leads to the most learning?
\end{itemize}
\end{frame}

%------------------------------------------------

\begin{frame}{Rubrics as Learning Tools}
We already use rubrics—but how?

\vspace{1em}

\textbf{Two possibilities:}
\begin{itemize}
\item Grading tool only
\item Learning tool (for students)
\end{itemize}

\vspace{1em}

\textbf{Growth-aligned use:}
\begin{itemize}
\item Students read rubrics before submission
\item Students use rubrics to self-assess
\item Rubrics guide revision
\end{itemize}
\end{frame}

%------------------------------------------------

\begin{frame}{Discussion: Rubrics in Our Courses}
\textbf{Consider:}
\begin{itemize}
\item Do students meaningfully engage with rubrics?
\item Are rubrics clear enough to guide improvement?
\item Could we simplify or restructure them?
\end{itemize}
\end{frame}

%------------------------------------------------

\begin{frame}{High-Leverage Small Changes}
We do \textbf{not} need a full redesign.

\vspace{1em}

\textbf{Possible small shifts:}
\begin{itemize}
\item Add one sentence of actionable feedback per problem
\item Delay release of scores (show feedback first)
\item Require one corrected solution per exam
\item Highlight “what was done well” explicitly
\end{itemize}
\end{frame}

%------------------------------------------------


%------------------------------------------------























%------------------------------------------------



%------------------------------------------------


%------------------------------------------------

%------------------------------------------------

%------------------------------------------------


%------------------------------------------------


%------------------------------------------------

%------------------------------------------------

\begin{frame}{Final Discussion}
\textbf{Actionable questions:}
\begin{itemize}
\item What is one small change we could make to exam feedback?
\item Where could we incorporate additional revision or reattempts?
\item How can we shift from grading to feedback (even partially)?
\end{itemize}

\vspace{1em}

\textbf{Goal:}
\begin{itemize}
\item Align our assessment practices with a growth mindset
\end{itemize}
\end{frame}

%------------------------------------------------

\end{document}


